\input{./relpath-to-latexroot.ltx}
\documentclass[\latexroot/\projectname]{subfiles}
\input{\latexroot/@resources/subfile-setup.ltx}

\begin{document}

% NOTE (2026-08-09, R-b/BUG-074 adoption): the live pipeline's pickle now
% carries ALL THREE policies under the fixed-real-rate rule with the tax
% cut household-financed (the published caption's original "under a fixed
% real rate rule" phrase becomes TRUE on the candidate figure). This
% caption stays regime-silent while the FROZEN published PDF (regime-mixed)
% is displayed; restore the fixed-real phrase WHEN the candidate is
% promoted (rides the R-f errata/corrections process).
\begin{figure}[htb]
  \centering
  \caption{\pregenmark Multiplier comparison: partial equilibrium vs.\ HANK model}
  \whenintegrated{\label{fig:HANK_multipliers}} % integrated doc: SST for labels
  \frozenORcandidate[width=.9\textwidth]{\latexroot/Code/HA-Models/FromPandemicCode/Figures/Cumulative_multipliers_withHank.pdf}

  \medskip
  \noindent\parbox{\textwidth}{\footnotesize
    \textbf{Note}: This figure compares consumption multipliers across model frameworks (Section~\ref{sec:hank}).
    The comparison shows multipliers over different time horizons
    between the baseline partial equilibrium model (with state-dependent multipliers during recessions)
    and the general equilibrium HANK-SAM model (with uniform Keynesian cross mechanism).
    Both models show the tax cut policy substantially underperforming relative to stimulus checks
    and UI extensions, validating the robustness of policy rankings across modeling approaches.
    The HANK model produces larger overall multipliers due to general equilibrium effects.
  }
\end{figure}

\vspace{0.5em}

% Smart bibliography: Only include bibliography if standalone AND has citations
\smartbib

\end{document}
